\documentclass[UTF8,a4paper]{ctexart}

\usepackage{geometry}
\usepackage{amsmath,amssymb,bm}
\usepackage{booktabs,longtable,array,tabularx}
\usepackage{graphicx,float,caption}
\usepackage{adjustbox,etoolbox}
\usepackage{fvextra}
\usepackage{fancyhdr}
\usepackage{xcolor}
\usepackage{placeins}
\usepackage[hidelinks]{hyperref}

\setmonofont{DejaVu Sans Mono}

% 页面参数仅取自用户指定参考稿 main-v4.1.tex 的页面设置。
\geometry{
  a4paper,
  top=25.4mm,
  bottom=25.4mm,
  left=31.75mm,
  right=31.75mm,
  headheight=15pt,
  headsep=10mm,
  footskip=12mm
}

% 取消页眉；所有页面只在页脚正中显示阿拉伯数字页码。
\pagestyle{fancy}
\fancyhf{}
\cfoot{\thepage}
\renewcommand{\headrulewidth}{0pt}
\renewcommand{\footrulewidth}{0pt}
\fancypagestyle{plain}{%
  \fancyhf{}%
  \cfoot{\thepage}%
  \renewcommand{\headrulewidth}{0pt}%
  \renewcommand{\footrulewidth}{0pt}%
}

% 一级标题使用中文序号；二、三级标题保持 1.1 与 1.1.1 的层级格式。
\ctexset{
  section={name={,、},number=\chinese{section}},
  subsection={name={,},number=\arabic{section}.\arabic{subsection}},
  subsubsection={name={,},number=\arabic{section}.\arabic{subsection}.\arabic{subsubsection}}
}

\graphicspath{{figures/}}
% 功能性溢出保护：不改变正文内容，只约束超宽表格并改善长英文词附近的断行。
\BeforeBeginEnvironment{tabular}{\begin{adjustbox}{max width=\linewidth}}
\AfterEndEnvironment{tabular}{\end{adjustbox}}
\setlength{\emergencystretch}{2em}

% 正式题名和队伍匿名信息须在官方规则确认后替换。
\title{海上人员运输的异构直升机路径与机队排班优化}
\author{}
\date{}


% ===== 首页版式修复：题名从正文版心顶部开始，避免默认 \maketitle 的额外下移 =====
\makeatletter
\renewcommand{\maketitle}{%
  \thispagestyle{plain}%
  \noindent\makebox[\textwidth][c]{\LARGE\bfseries\@title}\par
  \vspace{1.4em}%
}
\makeatother

% ===== 摘要版式修复：与正文同字号、同左右版心 =====
% ctexart 默认 abstract 会使用较小字号并缩进左右边界；此处取消该默认行为。
\renewenvironment{abstract}{%
  \par
  \noindent\makebox[\textwidth][c]{\bfseries 摘要}\par
  \vspace{0.75em}%
  \normalfont\normalsize
  \setlength{\parindent}{2em}%
  \setlength{\parskip}{0pt}%
  \ignorespaces
}{%
  \par
}

\begin{document}

\pagenumbering{arabic}
\setcounter{page}{1}
\maketitle

% mcm-paper source binding (non-rendering metadata)
% Repository HEAD: 66f49838fe3060c9c0286cb5aa87613ac20b9fc3
% Abstract content version: V0.4
% Recommended title: 海上人员运输的异构直升机路径与机队排班优化
% NOTE: formal filename remains 00_abstract.tex.

\begin{abstract}
海上设施人员运输同时受到异构直升机、座位与燃油、多机场闭环、海上着陆次数、人员取送、时间窗及多日机队排班等约束。针对题目三问，本文建立“统一架次可行性核—候选路线/时间列生成—整数主问题—词典序优化—独立验证”的递进式优化框架，在统一物理规则下完成单向出海、联合取送与有限机队多日调度。

针对问题一，构建有限物理路线池—人员流整数规划模型，将路线执行次数与请求组承运人数分离。候选路线由地理近邻设施簇、访问顺序枚举及燃油标号搜索生成，并结合解引导扩池与CP-SAT五级词典序优化完成求解。最终1600名人员全部完成运输，飞机使用时间为14759 min，共执行88架次，燃油消耗118693 kg。

针对问题二，构建物理路线池—服务选项动态人员流联合取送整数规划模型，以服务选项刻画不同人员在候选路线中的取送方式，并利用半开区间容量约束描述座位即时复用。通过确定性成对插入和解引导路线池扩展生成联合运输候选路线，再利用CP-SAT依次优化五级目标。最终4000名人员全部完成运输，飞机使用时间18963 min，共执行111架次，距离加权座位利用率84.95\%；相较独立OD运输基准，飞机使用时间、燃油消耗和架次数分别下降66.4\%、68.0\%和69.0\%，座位利用率提高57.5个百分点。

针对问题三，构建时间可行架次列—具体飞机指派—服务选项多日机队排班整数规划模型，将实际起飞时刻、任务时间窗、30 min周转及单机时序冲突纳入统一约束。根据任务时间窗和路线持续时间生成临界起飞时刻及时间可行架次列，并采用A--C三阶段词典序求解：先确定非临时任务基准，再在固定时间预算下最大化临时任务接纳人数，最后继续优化次级指标。160名临时人员全部接纳，其中阶段B在35114 min预算和当前有限时间列池内获得最优性认证；阶段C最终得到飞机使用时间34033 min、204架次、燃油消耗273629 kg，并启用22架飞机。

三问最终方案均通过独立校验程序复核人员完整性、容量、燃油、路线及相应时间约束；问题一的加油候选数量敏感性实验与问题三的时间列密度实验进一步检验了模型稳定性。结果表明，该框架能够有效协调路线共享、动态容量与机队时空资源，为复杂海上直升机运输计划提供可执行的优化方案。
\end{abstract}

\noindent\textbf{关键词：}海上人员运输；异构直升机；整数规划；词典序优化；机队排班

\clearpage
% mcm-paper source binding (non-rendering metadata)
% Repository HEAD: c9fab282b76ae99fb2878cbca036535471ea3123
% Section scope: problem restatement only; no model/algorithm/result claims added.

\section{问题重述}\label{sec:problem-restatement}

\subsection{问题背景}\label{subsec:problem-background}

海上油气设施分布广、陆海距离较远，人员往返及设施间调配主要依赖直升机完成。制定运载计划时，既要保证人员准确到达指定地点，又需满足机型容量、航程与燃油、停靠与加油、机场闭环等运行要求。题目给出了机场与海上设施距离、人员运输需求及直升机性能参数；单架次从陆地机场起飞，经若干海上设施完成运输后返回原机场，且访问设施数量受到限制。方案以总飞机使用时间为首要评价指标，同时考虑人员在途时间、燃油消耗、架次数和座位利用率。

随着约束逐步增加，问题由单向出海运输扩展为出海、海返与设施间穿梭的联合运输，进一步发展为具有人员时间要求和有限具体机队的多日排班。因此，需要在统一运行规则下依次确定人员分配、飞行路线及具体飞机排班方案。

\subsection{问题提出}\label{subsec:problem-statement}

\textbf{问题一：单向出海人员运载计划。}
对1600名出海人员，在暂不考虑具体时间要求且飞机数量不受限制的条件下，确定机型、起飞机场、飞行架次、海上设施访问顺序、人员搭载及必要的加油安排，使全部人员由陆地机场到达对应海上目的地。要求在满足统一飞行规则的前提下，以总飞机使用时间最小为首要目标，并兼顾其余评价指标，给出人员运输与飞行路线方案。

\textbf{问题二：多类型人员联合运载计划。}
将运输对象扩展至4000人，包括出海、海返和海上设施间穿梭三类需求。人员可在同一架次的不同停靠点依次上下机，使释放出的座位继续服务后续人员，因此需联合确定人员组合、停靠顺序、机型、机场及加油方案。在保证每名人员由同一架次完成运输的基础上，形成多方向需求协同的联合运载方案，并按统一指标评价运输效率。

\textbf{问题三：考虑时间要求与有限机队的多日运输计划。}
在7日计划周期内，使用题目指定的24架具体直升机完成具有不同时间要求的运输任务。各飞机具有固定基地和机型，并受每日起飞、返场及连续架次周转时间限制；人员运输还需满足最早出发与最晚到达要求。非临时任务必须完成，临时任务可根据剩余运输能力选择性接纳。首先以全部非临时任务建立飞机使用时间基准，再在不超过该基准的条件下尽可能接纳临时人员，并进一步优化其他指标，最终给出由具体飞机、执行架次、人员安排及各站到离时刻构成的7日完整排班方案。

\section{问题分析}\label{sec:problem-analysis}

三个问题共享机场、海上设施、机型、燃油及停靠等基本规则，但约束逐步增加：问题一处理单向出海运输，问题二扩展为多方向联合取送，问题三进一步加入时间窗、具体机队和多日排班。因此，三问可视为同一运输优化框架下逐层扩展的三个阶段。从车辆路径问题的视角看，异构车队、多中心与多架次均属于现实运输优化中常见的扩展结构，已有研究分别讨论了异构车队多中心协同配送以及异构车队多架次路径优化问题\cite{liang2024,sethanan2020}。

\subsection{问题一分析}\label{subsec:q1-analysis}

问题一包含1600名出海人员，运输方向均为陆地机场至海上设施。在没有海上上客的情况下，同一架次中的人员只会随停靠过程逐步下机，因此机上人数单调减少，最大载客量出现在起飞后的第一航段，这使问题一的容量关系相对简单。

该问题的关键在于，在满足机型容量、最多海上着陆次数、燃油安全及机场闭环等条件下，合理组合人员并确定海上设施访问顺序。不同路线不仅影响飞机总使用时间，也会影响人员在途时间、燃油消耗及座位利用率。因此，可将问题转化为可行飞行路线构造与人员分配相结合的运输优化问题，以总飞机使用时间为主要优化目标。

\subsection{问题二分析}\label{subsec:q2-analysis}

问题二将运输规模扩大至4000人，并新增海返人员和海上设施间穿梭人员。与问题一相比，主要变化不是人数增加，而是人员可以在飞行途中不同站点上下机，使飞机载客量随航段动态变化。

在同一设施完成下客后，释放出的座位可以立即用于后续人员上机，因此需要按照路线访问顺序逐航段计算实际载客人数。同时，海返和穿梭人员只有在飞机先访问其起点、再访问终点时才能完成运输，人员分配与路线顺序之间形成更强耦合。

因此，问题二应在问题一路线规划基础上进一步考虑人员取送顺序和动态座位复用，形成路线选择、人员上下机安排与航段容量联合优化的问题。

\subsection{问题三分析}\label{subsec:q3-analysis}

问题三在问题二的联合运输基础上进一步加入7日规划周期、24架具体直升机以及人员时间窗和任务类别，其中3840名非临时人员必须完成运输，160名临时人员允许选择性接纳。此时，一条空间上可行的运输路线还必须确定具体执行日期、起飞时刻和执行飞机，并满足人员最早出发、最晚到达以及飞机周转和每日运营时间等要求。

因此，问题三由前两问的路线与人员配置问题进一步扩展为人员、路线、时间与具体飞机的联合排班问题。已有海上油气人员直升机调度研究将同类问题描述为同时具有取送、异构机队、多基地、多架次和时间同步特征的复杂车辆路径问题\cite{nafstad2021}，这与本题第三问的主要约束结构具有较强对应性。同时，由于非临时任务必须优先保障，而临时任务允许选择性接纳，可先完成非临时任务并建立飞机使用时间基准，再在不超过该基准的条件下尽可能接纳临时人员；在固定临时人员接纳数量后，进一步按照既定优先级优化总飞机使用时间、人员在途时间、燃油消耗、架次数及座位利用率。

三个问题的核心递进关系可概括为
\[
\boxed{
\text{单向运输}
\rightarrow
\text{联合取送与动态容量}
\rightarrow
\text{时间窗约束下的有限机队排班}
}.
\]

\section{模型假设}\label{sec:assumptions}

为便于模型建立，在不改变题目基本运行规则的基础上，作如下假设：

\begin{enumerate}
    \item \textbf{忽略题目未给出的外部扰动及附加资源限制。}
    不考虑天气、风速、机械故障、临时空域变化等随机因素，同时不额外设置机场停机位、海上设施并发起降能力及加油设施服务能力等题目未规定的约束。

    \item \textbf{具有相同运输属性的人员具有等价性。}
    对起点、终点及其他相关属性完全一致的人员，可在求解阶段聚合为同一需求组，求解后再还原至具体人员；问题三还需保证组内人员的任务类型和时间要求一致。

    \item \textbf{问题三中的非临时任务均需完成。}
    将应急、增储上产和常规倒班任务视为必须满足的运输需求，仅临时任务允许根据剩余运输能力选择性安排。

    \item \textbf{问题三的时间要求通过时间可行架次列进行筛选。}
    候选方案仅保留满足人员最早上机、最晚到达及每日运行时段等时间规则的安排，不额外假设题目未规定的随机等待或延误；飞机返场后的30 min周转时间按题意计入具体飞机的连续占用约束。
\end{enumerate}

\section{符号说明}\label{sec:notation}

本文主要符号及含义如表~\ref{tab:notation}所示，其余仅在单个问题中使用的局部变量将在对应模型建立部分随式定义。

\begin{table}[htbp]
    \centering
    \caption{主要符号说明}
    \label{tab:notation}
    \begin{tabular}{c l c}
        \toprule
        符号 & 含义 & 单位或取值 \\
        \midrule
        $A$ & 陆地机场集合 & 3座 \\
        $F$ & 海上设施集合 & 52座 \\
        $K$ & 直升机机型集合 & T1、T2、T3 \\
        $G$ & 人员需求组集合 & -- \\
        $H$ & 问题三具体飞机集合 & 24架 \\
        $\mathcal L_r$ & 架次$r$的航段索引集合 & -- \\
        $d_{ij}$ & 地点$i$至地点$j$的飞行距离 & km \\
        $Q_k$ & 机型$k$的座位容量 & 人 \\
        $v_k$ & 机型$k$的飞行速度 & km/h \\
        $c_k$ & 机型$k$的单位距离油耗 & kg/km \\
        $q_g$ & 需求组$g$的人员数量 & 人 \\
        $T^{air}$ & 全部架次的总飞机使用时间 & min \\
        $T^{pax}$ & 全部人员的总在途时间 & 人$\cdot$min \\
        $F^{fuel}$ & 总燃油消耗 & kg \\
        $N^{flight}$ & 总架次数 & 次 \\
        $U$ & 总体座位利用率 & 无量纲 \\
        $B_0$ & 问题三采用的非临时任务时间基准 & min \\
        \bottomrule
    \end{tabular}
\end{table}

% mcm-paper source binding (non-rendering metadata)
% Repository HEAD checked: 66f49838fe3060c9c0286cb5aa87613ac20b9fc3
% Section role: shared data preprocessing and unified modeling foundation.
% 2026-08-17: only Figure 1 / Figure 2 image paths are remapped per user instruction.

\section{数据预处理与统一建模基础}\label{sec:data-foundation}

三问共享地点、机型和基本运行规则，仅需求类型、时间约束与机队限制逐步增加。为避免后文重复，本文先统一数据口径、架次表示、人员容量、燃油传播及评价指标，再分别建立三问模型。

\subsection{数据预处理}

\subsubsection{原始数据检查与标准化}

距离数据包含3座陆地机场和52座海上设施，共55个节点，其中8座海上设施允许加油。经检查，距离矩阵与人员文件中的地点编号一致，未发现缺失距离、非法地点或重复人员编号。问题一含1600名出海人员；问题二扩展为4000人，其中出海、海返和设施间穿梭分别为1600、1600和800人；问题三沿用4000人，并增加时间窗和任务类别。

\subsubsection{人员需求聚合与时间统一}

为压缩变量规模，本文将约束属性完全相同的人员聚合为需求组；原始人员编号仅用于求解后的逐人回填。问题一按起终点聚合为104组，问题二按起终点聚合为264组，问题三进一步区分时间窗和任务类别，共形成1877组。需求中的\texttt{LAND}标记表示可由满足条件的任一陆地机场完成运输，不能与明确指定A01、A02或A03的需求混同。

问题三以2026年8月3日00:00为时间原点，将全部时刻转为整数分钟参与求解，输出时再还原为实际日期和时刻。三问的数据规模及约束递进关系见图~\ref{fig:data-overview}。图中可见，问题一到问题二主要体现运输类型与人员规模的扩展，而问题三在人员总量不变的情况下，需求组数由264组增至1877组，说明时间窗与任务类别细分是后续时空调度规模增长的主要来源。

\begin{figure}[htbp]
    \centering
    % 用户指定映射：论文图1“运输网络与三问数据规模概览”
    % 实际文件名为 05_Fig1_三问统一建模框架与目标层级.png
    \includegraphics[width=0.96\textwidth]{figures/05_Fig1_三问统一建模框架与目标层级.png}
    \caption{运输网络与三问数据规模概览}
    \label{fig:data-overview}
\end{figure}

\subsection{统一建模基础}

三问均先生成满足公共运行规则的候选架次，再叠加相应人员流、时间和机队约束。问题一为单向出海运输；问题二引入双向取送与动态座位复用；问题三进一步加入带时刻候选列和具体飞机排班。统一框架如图~\ref{fig:unified-framework}所示。该图表明，三问并非彼此独立的模型，而是在同一物理可行性核上逐步增加人员流、时间与具体飞机约束，因此后文只需给出各问新增的模型结构。

\begin{figure}[htbp]
    \centering
    % 用户指定映射：论文图2“三问统一建模框架与目标层级”
    % 实际文件名为 05_Fig2_运输网络与三问数据规模概览.png
    \includegraphics[width=0.96\textwidth]{figures/05_Fig2_运输网络与三问数据规模概览.png}
    \caption{三问统一建模框架与目标层级}
    \label{fig:unified-framework}
\end{figure}

\subsubsection{架次表示与时间规则}

候选架次$r$由起降机场$b_r$、机型$k_r$、有序停靠序列及加油决策组成。若海上着陆次数为$m_r$，则地点序列满足

\[
v_{r,0}=v_{r,m_r+1}=b_r,
\qquad
1\le m_r\le5.
\tag{5-1}
\]

即每个架次从某一机场起飞，经至多5次海上着陆后返回同一机场。机型$k$从$i$到$j$的航段时间按题意逐段向上取整：

\[
\tau_{ij}^{k}
=
\left\lceil
\frac{60d_{ij}}{v_k}
\right\rceil .
\tag{5-2}
\]

海上设施基础停靠10 min，加油额外增加10 min；问题三再在物理路线基础上附加执行日期和完整到离时刻。

\subsubsection{人员取送与动态容量}

设人员在架次$r$上的上、下机位置为$\mu$和$\nu$，其占座区间定义为半开区间$[\mu,\nu)$。于是第$h$个航段的载客量为

\[
n_{r,h}
=
\sum_{o\in\mathcal O_r}
z_o\mathbf 1(\mu_o\le h<\nu_o),
\qquad
n_{r,h}\le Q_{k_r}.
\tag{5-3}
\]

半开区间使同一站点下机人员释放的座位可立即供新人员使用。问题一只有机场上客，容量约束可进一步化简；问题二和问题三保留逐航段动态容量。

\subsubsection{燃油传播与加油规则}

所有架次从机场满油起飞。若机型$k_r$单位距离耗油为$c_{k_r}$，则逐航段燃油传播满足

\[
f_{r,h+1}^{arr}
=
f_{r,h}^{dep}
-
c_{k_r}d_{v_{r,h},v_{r,h+1}},
\qquad
f_{r,h+1}^{arr}\ge R_{k_r}.
\tag{5-4}
\]

仅8座指定海上设施允许加油；发生加油时，离站前燃油恢复至该机型满油状态。燃油受限车辆路径研究通常也需要同时处理续航约束、可加油节点以及从基地出发并返回基地的闭环路线\cite{sundar2017}。因此，本文在候选路线生成阶段对每一航段进行燃油传播与安全余油检查，并同步记录合法加油决策。

\subsubsection{统一评价指标}

三问统一记录总飞机使用时间$T^{air}$、人员总在途时间$T^{pax}$、燃油消耗$F^{fuel}$、架次数$N^{flight}$和距离加权座位利用率$U$。其中

\[
T^{air}
=
\sum_{r\in S}T_r,
\qquad
U
=
\frac{\sum_{r\in S}PK_r}
{\sum_{r\in S}AK_r}.
\tag{5-5}
\]

问题一、问题二及问题三的阶段C均按
$T^{air}\rightarrow T^{pax}\rightarrow F^{fuel}\rightarrow N^{flight}\rightarrow -U$
逐级优化；问题三的阶段B则在固定飞机时间预算下优先处理临时任务接纳。后文各模型均在上述公共规则上增加本问特有约束，不再重复定义。

% mcm-paper source binding (non-rendering metadata)
% Repository HEAD: 66f49838fe3060c9c0286cb5aa87613ac20b9fc3
% Q1 model reconciliation: 3919ddf34054d88f6f6b262b9a91967397b055d2
% Q1 formal result source: c9ae677177845579364c9b245554db79a3492d4a
% Q1 figure source: 6f0aae699bb2c17d5a58532699313bb8346c8861

\section{问题一模型的建立与求解}\label{sec:q1}

问题一仅含由陆地机场至海上设施的单向运输，无海上上客、人员时间窗和具体飞机排班。基于第~\ref{sec:data-foundation}节的统一可行性核，本文建立有限物理路线池--人员流整数规划模型：候选路线只确定机场、机型、停靠顺序和加油方案，人员配置由整数流变量在主问题中统一决定。

\subsection{有限物理路线池--人员流整数规划模型建立}\label{subsec:q1-model}

\subsubsection{决策变量与需求平衡约束}

记需求组集合为$G_1$，组$g$人数为$q_g$；当前有限合法物理路线池为$\Omega$，可服务组$g$的路线集合为$\Omega(g)$。定义
\[
y_r\in\mathbb Z_+,\qquad
z_{g,r}\in\mathbb Z_+,
\]
其中$y_r$表示路线$r$的执行架次数，$z_{g,r}$表示需求组$g$由路线$r$承运的人数。

问题一所有人员均在机场上机、途中只下不上，因此任一路线的载客量随航段单调不增，最大值出现在首航段。故逐航段容量约束可等价简化为首航段的聚合容量约束；这一性质仅适用于问题一。

\subsubsection{目标函数与优化层级}

记路线$r$的飞机使用时间、燃油消耗和可用座公里分别为$T_r,F_r,AK_r$，组$g$由路线$r$运输时的在途时间和客公里为$A_{g,r},K_{g,r}$。五项指标写为

\[
\begin{aligned}
T_1^{air}&=\sum_{r\in\Omega}T_ry_r, &
T_1^{pax}&=\sum_{g}\sum_{r\in\Omega(g)}A_{g,r}z_{g,r},\\
F_1^{fuel}&=\sum_{r\in\Omega}F_ry_r, &
N_1^{flight}&=\sum_{r\in\Omega}y_r,\\
U_1&=
\frac{\sum_g\sum_{r\in\Omega(g)}K_{g,r}z_{g,r}}
{\sum_{r\in\Omega}AK_ry_r}.
\end{aligned}
\tag{6-1}\label{eq:q1-metrics}
\]

令$\mathbf{x}=(y,z)$，并按
$T_1^{air}\rightarrow T_1^{pax}\rightarrow F_1^{fuel}\rightarrow N_1^{flight}\rightarrow -U_1$
进行词典序优化，则问题一正式模型统一写为

\[
\begin{array}{@{}l@{\qquad}c@{}}
\displaystyle \min_{\mathbf{x}}
&
\displaystyle
f_1(\mathbf{x})
=
\operatorname{lex}\!\left(
T_1^{air},T_1^{pax},F_1^{fuel},N_1^{flight},-U_1
\right)
\\[1.0em]
\mathrm{s.t.}
&
\left\{
\begin{aligned}
\sum_{r\in\Omega(g)} z_{g,r} &= q_g,
&& g\in G_1,\\
\sum_{g:\,r\in\Omega(g)} z_{g,r} &\le Q_{k_r}y_r,
&& r\in\Omega,\\
y_r &\in\mathbb Z_+,
&& r\in\Omega,\\
z_{g,r} &\in\mathbb Z_+,
&& g\in G_1,\ r\in\Omega(g).
\end{aligned}
\right.
\end{array}
\tag{6-2}\label{eq:q1-model}
\]

式~\eqref{eq:q1-model}中第一条保证需求完整运输，第二条为问题一特有的聚合容量约束。更高优先级目标取得当前最好值后固定，再继续优化下一目标；分式利用率采用Dinkelbach迭代处理。

\subsection{问题一模型求解}\label{subsec:q1-solution}

\subsubsection{候选物理路线生成}

完全枚举52座海上设施的1--5元组合规模过大，因此首先采用“种子设施+地理近邻”的方式生成候选设施簇。对每个种子设施，仅从其若干最近邻中组合其余成员，从而将候选规模控制在可求解范围内；该步骤属于启发式构池，也是当前有限路线池与完整可行空间之间的主要差别之一。

固定机场$b$、机型$k$及候选设施簇$S$后，对$m=|S|\le5$个必访设施的访问顺序进行完整枚举，最多需要考察$5!=120$种排列。对于每一种访问顺序，内层采用Dijkstra式燃油标号搜索决定是否插入额外加油停靠以及具体加油位置。搜索状态可表示为
\[
(v,\,p,\,n_{\mathrm{ref}}),
\]
其中$v$为当前位置，$p$为已经完成的必访设施数，$n_{\mathrm{ref}}$为已插入的额外加油停靠数；每个状态对应的标号记录累计飞机使用时间和当前剩余燃油。由当前位置向下一必访设施或候选加油设施扩展时，按航段距离扣除燃油并累加逐段向上取整后的飞行时间；若到达后的剩余燃油低于安全余油，或累计海上着陆次数超过5次，则直接舍弃该标号。若在合法加油设施执行加油，则额外计入加油停靠时间，并将离站燃油恢复至该机型满油状态。由此可在给定访问顺序和加油候选集合内搜索得到飞机使用时间最小的合法加油方案。

为控制内层搜索规模，正式运行并未将8座合法加油设施全部作为额外插入候选，而是默认保留与当前机场和任务设施最相关的4座加油设施。因此，上述标号搜索应理解为在给定必访设施集合、访问顺序及所选加油候选子集内的精确搜索，而不等同于完整路线空间上的全局最优搜索。所有生成路线随后通过第~\ref{sec:data-foundation}节的统一物理可行性检查和支配筛选后进入主问题。正式运行由3200个候选设施簇生成31257条物理路线，并进一步筛选得到初始主问题路线池。

\subsubsection{整数主问题与解引导扩池}

首先用贪心方案为CP-SAT提供初始可行解，再求解式~\eqref{eq:q1-model}。随后依据当前已选路线的设施组合构造子集、并集和单点替换邻域，重新生成有潜力的路线并回灌主问题。正式运行进行了4轮解引导扩池，路线池由2220条扩展至2490条。如图~\ref{fig:q1-process}所示，扩池第1轮至第4轮对应4次定向候选池补充；每轮均依据当前已选路线构造邻域并重新生成候选路线，以逐步扩大高质量组合空间。

主目标稳定后依次进行人员在途时间、燃油、架次数和利用率优化。最终聚合人员流被拆分为具体架次并按原始人员编号回填，再由独立校验程序重新检查人员完整性、容量、燃油和全部评价指标。

\subsection{问题一求解结果及分析}\label{subsec:q1-results}

\subsubsection{总体结果与求解过程}

正式求解得到当前有限路线池下的最好已验证可行方案，主要指标汇总于表~\ref{tab:q1-summary}。

\begin{table}[htbp]
\centering
\caption{问题一当前正式方案主要指标}
\label{tab:q1-summary}
\small
\begin{tabular}{c c c c c}
\toprule
$T_1^{air}$/min & $T_1^{pax}$/人$\cdot$min & 燃油/kg & 架次数 & $U_1$\\
\midrule
14759 & 126288 & 118693.0 & 88 & 49.64\%\\
\bottomrule
\end{tabular}
\end{table}

如图~\ref{fig:q1-process}所示，正式求解链条从贪心可行解16318 min出发，经主问题初次求解降至14886 min，并在4轮解引导扩池后进一步降至14759 min。主目标在后续扩池中仍持续改善，说明围绕当前解补充候选路线能够有效扩大高质量组合空间。LP松弛参考下界为14247.6 min，当前CP-SAT最好界为14470 min，因此

\[
14470\le T_{1,\Omega}^{*}\le\widehat T_{1,\Omega}=14759,
\tag{6-3}\label{eq:q1-bound}
\]

当前可行值与求解器界相差1.96\%。结合图~\ref{fig:q1-process}中的单调改进过程可知，当前路线池已经获得较稳定的高质量解，但尚不能据此宣称池内或原问题全局最优。

\begin{figure}[htbp]
    \centering
    \includegraphics[width=0.90\textwidth]{figures/04求解过程及飞机使用时间变化.pdf}
    \caption{问题一求解过程及飞机使用时间变化}
    \label{fig:q1-process}
\end{figure}

图~\ref{fig:q1-lex-status}给出了五级目标的求解状态、相对差距与计算代价。各阶段均获得可行解但未完成完整词典序最优性认证，其中人员在途时间阶段的相对差距约为8.73\%，明显高于主目标和燃油阶段，是当前词典序链中最主要的未收敛环节；利用率采用Dinkelbach迭代处理，其计算时间也明显高于前四级目标。因此，最终方案应理解为在既定优先级下逐级改进得到的经验证可行方案。

\begin{figure}[htbp]
    \centering
    \includegraphics[width=0.90\textwidth]{figures/05五级目标的求解状态与相对差距.pdf}
    \caption{问题一五级目标的求解状态与相对差距}
    \label{fig:q1-lex-status}
\end{figure}

\subsubsection{运输方案结构分析}

最终88个架次中，较大容量的T2和T3承担全部运输任务，其中83个架次满载起飞，占全部架次的94.3\%；全部架次首航段共提供1609个座位，其中1600个被占用，因此整体起飞座位装载率达到99.44\%。但总体距离加权座位利用率仅为49.64\%，原因是人员沿多设施路线逐站下机，返场及后续航段仍产生可用座公里。

如图~\ref{fig:q1-leg-utilization}所示，真实T3-48架次对此给出直观解释：飞机由A03满载19人起飞，经F036、F031、F023和F046后返场，并在F031完成加油；机上人数随各站下客由19依次降至15、10、1和0，因此该架次虽然满载起飞，距离加权利用率仍为47.25\%。这既验证了问题一中聚合容量约束的合理性，也说明“起飞满载”与“全航程高利用率”并不等价；提高总体座位利用率的关键不仅是首航段装满，还在于减少后续低载客航段和空返距离。

\begin{figure}[htbp]
    \centering
    \includegraphics[width=0.92\textwidth]{figures/06逐航段载客量与距离加权座位利用率.pdf}
    \caption{问题一逐航段载客量与距离加权座位利用率}
    \label{fig:q1-leg-utilization}
\end{figure}

综上，有限物理路线池与人员流联合优化能够完整运输1600名人员，并通过解引导扩池持续改善主目标；最终方案已通过独立验证。由于候选池采用启发式构造且各词典序阶段均只获得可行解、尚未完成最优性证明，14759 min仅作为当前有限路线池下的最好已验证可行值。

% mcm-paper source binding (non-rendering metadata)
% Repository HEAD: 66f49838fe3060c9c0286cb5aa87613ac20b9fc3
% Q2 model reconciliation: 3919ddf34054d88f6f6b262b9a91967397b055d2
% Q2 formal result source: c9ae677177845579364c9b245554db79a3492d4a
% Q2 figure source: 5def82abd520b763732631fa1474ab83c202a273

\section{问题二模型的建立与求解}\label{sec:q2}

问题二同时包含出海、海返和设施间穿梭需求，人员可在同一架次不同停靠点上下机。与问题一相比，核心变化是座位需求不再沿航程单调下降，因此必须显式描述人员在各航段的占座区间。取送类车辆路径问题的相关研究表明，取送先后关系、车辆容量以及时间或同步条件往往需要联合考虑\cite{ghilas2016alns}。结合本题的多类运输需求，本文进一步建立有限物理路线池与动态人员流联合取送模型。

\subsection{物理路线池--动态人员流联合取送模型建立}\label{subsec:q2-model}

\subsubsection{服务选项与动态容量}

记需求组集合为$G_2$，有限物理路线池为$\Omega$。对需求组$g$和路线$r$，若路线依次经过其起点和终点，则构造服务选项
$o=(g,r,\mu_o,\nu_o)$，其中$\mu_o$、$\nu_o$分别为上、下机位置，且终点取上机后首次到达的位置。定义

\[
y_r\in\mathbb Z_+,\qquad z_o\in\mathbb Z_+,
\]

分别表示路线$r$的执行架次数和服务选项$o$承运的人数。人员仅在半开区间$[\mu_o,\nu_o)$内占座，因此路线$r$第$h$个航段的载客量为

\[
n_{r,h}
=
\sum_{\substack{o\in\mathcal O_r\\ \mu_o\le h<\nu_o}}z_o.
\tag{7-1}\label{eq:q2-leg-load}
\]

半开区间使同一停靠点下机释放的座位可立即供新人员使用。如图~\ref{fig:q2-dynamic-capacity}所示，真实T2-21架次累计运输59人，而机型容量仅16人；由于人员在不同停靠点先后下机、上机，各航段最大在机人数始终不超过16。该例说明，问题二的容量可行性应由逐航段同时占座人数判断，而不能用单架次累计服务人数直接与机型容量比较。

\begin{figure}[htbp]
    \centering
    \includegraphics[width=0.94\textwidth]{figures/10半开区间动态容量与座位即时复用.pdf}
    \caption{问题二半开区间动态容量与座位即时复用}
    \label{fig:q2-dynamic-capacity}
\end{figure}

\subsubsection{正式优化模型}

记服务选项$o$的人员在途时间和客公里为$P_o,K_o$，路线$r$的飞机使用时间、燃油和可用座公里为$T_r,F_r,AK_r$。五项评价指标为

\[
\begin{aligned}
T_2^{air}&=\sum_{r\in\Omega}T_ry_r, &
T_2^{pax}&=\sum_{o\in\mathcal O}P_oz_o,\\
F_2^{fuel}&=\sum_{r\in\Omega}F_ry_r, &
N_2^{flight}&=\sum_{r\in\Omega}y_r,\\
U_2&=\frac{\sum_{o\in\mathcal O}K_oz_o}
{\sum_{r\in\Omega}AK_ry_r}.
\end{aligned}
\tag{7-2}\label{eq:q2-metrics}
\]

令$\mathbf{x}=(y,z)$，问题二按统一五级词典序求解，其正式模型写为

\[
\begin{array}{@{}l@{\qquad}c@{}}
\displaystyle \min_{\mathbf{x}}
&
\displaystyle
f_2(\mathbf{x})
=
\operatorname{lex}\!\left(
T_2^{air},T_2^{pax},F_2^{fuel},N_2^{flight},-U_2
\right)
\\[1.0em]
\mathrm{s.t.}
&
\left\{
\begin{aligned}
\sum_{o\in\mathcal O_g}z_o &= q_g,
&& g\in G_2,\\
\sum_{\substack{o\in\mathcal O_r\\ \mu_o\le h<\nu_o}}z_o
&\le Q_{k_r}y_r,
&& r\in\Omega,\ h\in\mathcal L_r,\\
\sum_{o\in\mathcal O_r}z_o=0
&\Rightarrow y_r=0,
&& r\in\Omega,\\
y_r&\in\mathbb Z_+,
&& r\in\Omega,\\
z_o&\in\mathbb Z_+,
&& o\in\mathcal O.
\end{aligned}
\right.
\end{array}
\tag{7-3}\label{eq:q2-model}
\]

式~\eqref{eq:q2-model}中，第一条保证4000名人员均被准确运输，第二条为逐航段动态容量约束，第三条排除无人员流的空架次。更高优先级目标固定当前取得值后再进入下一阶段，利用率仍采用Dinkelbach迭代。

\subsection{问题二模型求解}\label{subsec:q2-solution}

\subsubsection{候选路线池构造}

初始路线池由确定性规则生成：除每个OD的单组可行路线外，还加入共享起点、共享终点、前一需求终点与后一需求起点衔接以及出海--海返成对插入路线。固定访问顺序后，通过燃油标号同时决定合法加油位置，并使用全部8座可加油设施；所有路线均再次通过第~\ref{sec:data-foundation}节的统一可行性检查。

在当前解基础上继续进行解引导路线池扩展，定向补充有潜力的联合运输路线并重新求解。具体采用确定性成对插入与基于当前解的定向扩池方式构造新的联合运输候选。正式主问题包含10565条物理路线、87977个服务选项和44330组动态容量约束。

\subsubsection{词典序求解与回填}

整数主问题采用CP-SAT按既定优先级逐级求解。整数人员流求得后，将$[\mu_o,\nu_o)$区间展开并分配至具体架次，从而形成可执行运输方案；同站下机与上机可在同一位置完成座位复用。最终4000名人员按原始编号回填，并由独立校验程序重新检查需求、逐航段容量、燃油和全部指标。

如图~\ref{fig:q2-lex-status}所示，五级目标均获得可行解，但尚未完成完整最优性证明。主目标和人员在途时间的当前可行值与求解界差距分别为28.6\%和7.1\%，表明主目标仍存在较大的认证空间；后续目标又建立在未认证的上一级当前解上，因此图中各阶段反映的是逐级获得并稳定改进可行方案的计算过程，而非完整词典序最优性证明。

\begin{figure}[htbp]
    \centering
    \includegraphics[width=0.84\textwidth]{figures/09五级词典序求解状态与计算代价.pdf}
    \caption{问题二五级词典序求解状态与计算代价}
    \label{fig:q2-lex-status}
\end{figure}

\subsection{问题二求解结果及分析}\label{subsec:q2-results}

\subsubsection{联合运输总体效果}

当前有限路线池下得到的最好已验证可行方案为

\[
T_2^{air}=18963\ \mathrm{min},\quad
T_2^{pax}=241751\ \text{人}\!\cdot\!\mathrm{min},
\]

总燃油150491.1 kg，共111个架次，总体座位利用率84.95\%。主目标求解器下界为13548 min，故

\[
13548\le T_{2,\Omega}^{*}
\le \widehat T_{2,\Omega}=18963.
\tag{7-4}\label{eq:q2-bound}
\]

如图~\ref{fig:q2-baseline}所示，与独立OD运输基准相比，联合方案将飞机使用时间、燃油和架次数分别降低66.4\%、68.0\%和69.0\%，座位利用率提高57.5个百分点；代价是人员总在途时间增加4.5\%。架次数大幅下降与座位利用率显著提高同时出现，说明联合运输的主要收益来自跨OD需求共享同一物理架次和航段座位，而不是简单缩短某一条路线。这里的独立OD方案仅作为可复现对照，不是理论下界或“分离运输最优方案”。

\begin{figure}[H]
    \centering
    \includegraphics[width=0.84\textwidth]{figures/07Q2联合运输与独立OD运输基线对比.pdf}
    \caption{问题二联合运输与独立OD运输基线对比}
    \label{fig:q2-baseline}
\end{figure}

\subsubsection{人员时间的结构性权衡}

图~\ref{fig:q2-pax-tradeoff}进一步表明，人员时间增加并非三类需求普遍恶化。联合运输后，出海和海返人员平均在途时间分别由69.77、67.39 min升至71.44、72.46 min，而设施间穿梭由14.92 min降至14.39 min；对应总在途时间变化为$+2665$、$+8123$和$-426$人$\cdot$min。约16.6\%的人员发生额外经停，主要来自出海与海返需求为提高航段复用而合并运输。由此可见，人员时间成本的增加集中在少数合并运输路径，而穿梭任务反而因航段复用获得一定时间收益。

\begin{figure}[H]
    \centering
    \includegraphics[width=0.84\textwidth]{figures/08人员在途时间的结构性权衡.pdf}
    \caption{问题二人员在途时间的结构性权衡}
    \label{fig:q2-pax-tradeoff}
\end{figure}

因此，问题二的主要收益并非“五项指标同时改善”，而是利用跨OD合乘和座位即时复用显著减少飞机运行资源，以少量人员时间成本换取更高的整体运营效率。最终联合方案全部通过独立验证，但由于有限路线池及各阶段均未获得完整最优性认证，18963 min仍作为当前最好已验证可行值。

\FloatBarrier
% mcm-paper source binding (non-rendering metadata)
% Repository HEAD: 66f49838fe3060c9c0286cb5aa87613ac20b9fc3
% Q3 model reconciliation: 3919ddf34054d88f6f6b262b9a91967397b055d2
% Q3 formal 阶段B/C update source: 6f0aae699bb2c17d5a58532699313bb8346c8861
% Q3 figure source: a3970ef9eb51b87bd1c564a98d2cdfef213373a4

\section{问题三模型的建立与求解}\label{sec:q3}

问题三在联合取送基础上进一步加入7日计划周期、24架具体飞机、人员时间窗及30 min周转约束。此时路线不仅要在空间上可行，还必须确定执行日期、起飞时刻和具体飞机。本文建立时间可行架次池--具体飞机排班与人员流选择性服务模型，并按“非临时任务基准--临时任务接纳--次级指标细化”三阶段求解。

\subsection{时间可行架次池--具体飞机排班模型建立}\label{subsec:q3-model}

\subsubsection{带时刻候选列与公共可行域}

记非临时、临时需求组分别为$G_3^M,G_3^T$。一条带时刻候选列$r$固定日期、基地$b_r$、机型$k_r$、停靠与加油方案、起飞时刻$s_r$及返场时刻$e_r$；返场后30 min计入具体飞机占用区间：

\[
I_r=[s_r,e_r+30).
\tag{8-1}\label{eq:q3-occupancy-interval}
\]

对与列$r$基地和机型兼容的具体飞机$\alpha$定义$x_{r,\alpha}\in\{0,1\}$，并令
$y_r=\sum_{\alpha\in H_r^{comp}}x_{r,\alpha}$。
服务选项$o=(g,r,\mu_o,\nu_o)$表示需求组$g$在列$r$的$\mu_o$处上机、在首次到达终点的$\nu_o$处下机，变量$z_o\in\mathbb Z_+$表示承运人数。只有满足人员最早上机、最晚到达以及机场兼容条件的选项才进入候选集。

对具体飞机$\alpha$，定义时间冲突列对集合
\[
\mathcal C_\alpha=
\left\{
(r,r'):
r\neq r',\ 
\alpha\in H_r^{comp}\cap H_{r'}^{comp},\
I_r\cap I_{r'}\neq\varnothing
\right\}.
\tag{8-2}\label{eq:q3-conflict-set}
\]

令$\mathbf{x}=(x,y,z)$，三阶段共享的物理与排班可行域$\mathcal F_3^0$由

\[
\mathcal F_3^0=
\left\{
\begin{aligned}
y_r&=\sum_{\alpha\in H_r^{comp}}x_{r,\alpha},
&&r\in\Omega_3,\\
\sum_{\substack{o\in\mathcal O_r\\ \mu_o\le h<\nu_o}}z_o
&\le Q_{k_r}y_r,
&&r\in\Omega_3,\ h\in\mathcal L_r,\\
x_{r,\alpha}+x_{r',\alpha}&\le1,
&&(r,r')\in\mathcal C_\alpha,\ \alpha\in H,\\
\sum_{o\in\mathcal O_r}z_o=0&\Rightarrow y_r=0,
&&r\in\Omega_3,\\
x_{r,\alpha}&\in\{0,1\},\quad z_o\in\mathbb Z_+.
\end{aligned}
\right.
\tag{8-3}\label{eq:q3-feasible-set}
\]

其中动态容量仍采用半开区间$[\mu_o,\nu_o)$；每日06:00--18:00起飞、20:00前返场、同日运行及燃油可行性均已在时间可行架次列生成阶段保证。

\subsubsection{阶段A：非临时任务基准}

阶段A仅服务3840名非临时人员，以总飞机使用时间为目标：

\[
\begin{array}{@{}l@{\qquad}c@{}}
\displaystyle \min_{\mathbf{x}}
&
\displaystyle
f_A(\mathbf{x})
=
T_A^{air}
=
\sum_{r\in\Omega_3}\sum_{\alpha\in H_r^{comp}}T_rx_{r,\alpha}
\\[1.0em]
\mathrm{s.t.}
&
\left\{
\begin{aligned}
\mathbf{x}&\in\mathcal F_3^0,\\
\sum_{o\in\mathcal O_g}z_o&=q_g,
&&g\in G_3^M,\\
z_o&=0,
&&o\in\mathcal O_g,\ g\in G_3^T.
\end{aligned}
\right.
\end{array}
\tag{8-4}\label{eq:q3-stage-a}
\]

正式求解完成全部3840名非临时人员运输，得到当前最好已验证基准

\[
\widehat T_{0,\Omega}=35114\ \mathrm{min},
\qquad
25104\le T_{0,\Omega}^{*}\le35114.
\tag{8-5}\label{eq:q3-stage-a-bound}
\]

由于阶段A仅获得经验证的可行解、尚未完成最优性证明，35114 min作为后续阶段的实际预算基准，而不表述为全局最优值。记

\[
B_0=\widehat T_{0,\Omega}=35114\ \mathrm{min}.
\tag{8-6}\label{eq:q3-budget}
\]

\subsubsection{阶段B：临时任务接纳}

阶段B重新允许160名临时人员进入模型，在不突破$B_0$的前提下最大化接纳人数。为统一采用最小化形式，写为

\[
\begin{array}{@{}l@{\qquad}c@{}}
\displaystyle \min_{\mathbf{x}}
&
\displaystyle
f_B(\mathbf{x})
=
-\sum_{g\in G_3^T}\sum_{o\in\mathcal O_g}z_o
\\[1.0em]
\mathrm{s.t.}
&
\left\{
\begin{aligned}
\mathbf{x}&\in\mathcal F_3^0,\\
\sum_{o\in\mathcal O_g}z_o&=q_g,
&&g\in G_3^M,\\
0\le\sum_{o\in\mathcal O_g}z_o&\le q_g,
&&g\in G_3^T,\\
T^{air}(\mathbf{x})&\le B_0.
\end{aligned}
\right.
\end{array}
\tag{8-7}\label{eq:q3-stage-b}
\]

最新正式运行得到160/160名临时人员全部接纳，且当前最好可行值与求解器最优界均为160、相对间隙为0，故在固定$B_0$和当前有限时间可行架次列池下

\[
N_{\mathrm{temp},\Omega}^{*}=160.
\tag{8-8}\label{eq:q3-temp-optimum}
\]

对应总飞机使用时间为34668 min，仍比预算保留446 min余量。

\subsubsection{阶段C：固定服务水平后的词典序细化}

阶段C固定4000名人员全部完成运输，并按统一五级指标继续优化。记
$f_C(\mathbf{x})=\operatorname{lex}
(T_3^{air},T_3^{pax},F_3^{fuel},N_3^{flight},-U_3)$，则

\[
\begin{array}{@{}l@{\qquad}c@{}}
\displaystyle \min_{\mathbf{x}}
&
\displaystyle f_C(\mathbf{x})
\\[1.0em]
\mathrm{s.t.}
&
\left\{
\begin{aligned}
\mathbf{x}&\in\mathcal F_3^0,\\
\sum_{o\in\mathcal O_g}z_o&=q_g,
&&g\in G_3^M,\\
\sum_{o\in\mathcal O_g}z_o&=q_g,
&&g\in G_3^T,\\
T^{air}(\mathbf{x})&\le B_0.
\end{aligned}
\right.
\end{array}
\tag{8-9}\label{eq:q3-stage-c}
\]

前一级取得的当前目标值固定后再进入下一阶段。阶段C的前四级目标均获得可行解，第五级座位利用率优化在共享时间预算内未得到新的有效改进，因此最终正式方案采用第四级架次数优化所得的当前最好可行解，而不宣称获得完整五级词典序最优解。

\subsection{问题三模型求解}\label{subsec:q3-solution}

\subsubsection{时间可行架次池构造}

若逐分钟枚举路线、人员和具体飞机，变量规模将不可接受。海上直升机调度的相关研究曾采用预先生成可行架次、再在扩展模型中进行组合优化的思路，以降低直接刻画完整时空决策的难度\cite{nafstad2021}。与此类似，本文首先从问题二已采用的联合路线及非临时单OD短路线构造178条物理骨架，再将人员时间窗转化为每条路线的可行起飞区间。对每个非空区间仅取最早、中点和最晚三个临界时刻，并与每日运行时段相交，最终生成14634条时间可行架次列和90799个合法服务选项。

在有限列池上采用CP-SAT统一决定带时刻架次、具体飞机和人员流；单机冲突在数学模型中由式~\eqref{eq:q3-feasible-set}的冲突不等式刻画，求解时采用等价的不重叠全局约束提高调度约束的处理效率。阶段A、B、C依次运行，前一阶段的基准或服务水平进入下一阶段；聚合解最终回填到具体人员和飞机，并由独立校验程序重新检查时间窗、容量、燃油、30 min周转和全部评价指标。

\subsection{问题三求解结果及分析}\label{subsec:q3-results}

\subsubsection{三阶段递进结果}

图~\ref{fig:q3-stage-pipeline}集中展示三阶段结果。阶段A以35114 min完成3840名非临时人员；阶段B在不增加该基准的条件下将160名临时人员全部吸收；阶段C固定4000人全部服务后又将总飞机使用时间降至34033 min，并把架次数由阶段A的212次压缩到204次。图底部的独立验证及求解结果—交付方案一致性核对均通过，说明三阶段优化结果不仅在求解器内部可行，而且与最终人员和架次交付文件保持一致。

\begin{figure}[htbp]
    \centering
    \includegraphics[width=0.96\textwidth]{figures/11三阶段递进.pdf}
    \caption{问题三三阶段递进结果}
    \label{fig:q3-stage-pipeline}
\end{figure}

最终正式方案为

\[
T_3^{air}=34033\ \mathrm{min},\quad
T_3^{pax}=246808\ \text{人}\!\cdot\!\mathrm{min},
\]

总燃油273629.0 kg，共204个架次，使用22架具体飞机，总体座位利用率47.77\%。由于阶段C并未完成第五级认证，这组数值应理解为当前正式交付可行方案。

\subsubsection{具体飞机排班可执行性}

图~\ref{fig:q3-gantt}给出了24架飞机在8月3日至9日的完整甘特排班。204个架次均分配到具体飞机，实际投入22架；同一飞机的任务区间不存在重叠，最小实际周转时间恰为30 min。图右侧可见T2承担的架次数与累计飞机时间最多，T3次之，T1任务量较小，体现了容量、基地兼容和时间窗共同作用下的非均匀机队使用；同时22/24的实际投入规模说明模型保留了少量机队冗余，而不是依赖全部飞机满负荷运行。

\begin{figure}[htbp]
    \centering
    \includegraphics[width=\textwidth]{figures/12七天机队排班甘特图.pdf}
    \caption{问题三七天机队排班甘特图}
    \label{fig:q3-gantt}
\end{figure}

\subsubsection{临时任务全接纳与兼容性}

图~\ref{fig:q3-temp-compatibility}进一步解释160名临时人员为何能够全部纳入。所有临时任务时间窗均较宽，多数人员可兼容2--5架具体飞机，并具有一定的机型或基地替代空间；其中T2可服务99人，覆盖范围最大，T3和T1分别为62人和27人。需要注意，同一人员可能兼容多个机型或基地，因此这些数量不是互斥分类。

在$B_0=35114$ min下，阶段B实际仅使用34668 min，保留446 min飞机时间余量；同时临时人员接纳人数的当前最好可行值与最优界均为160，相对间隙为0，求解器完成该目标的最优性证明。因此，临时任务全部接纳并非依赖额外增加阶段A预算，而是通过联合排班与资源重组实现；兼容性分布同时说明，多机型、多基地替代空间是吸收临时需求的重要调度余度。

\begin{figure}[htbp]
    \centering
    \includegraphics[width=0.96\textwidth]{figures/13临时任务全接纳与兼容性画像.pdf}
    \caption{问题三临时任务全接纳与兼容性画像}
    \label{fig:q3-temp-compatibility}
\end{figure}

综上，问题三通过时间可行架次列、具体飞机区间排班和动态人员流，将问题二扩展到时间窗与有限机队环境。当前方案完成全部4000人运输，并通过独立验证；其中阶段B的临时人员接纳目标已在固定$B_0$和当前有限列池下严格认证，而阶段A及阶段C的完整最优性仍不作扩大解释。

% mcm-paper source binding (non-rendering metadata)
% Repository HEAD: 66f49838fe3060c9c0286cb5aa87613ac20b9fc3
% Validation report source: 代码/模型检验与最终结果报告.md @ a3970ef9eb51b87bd1c564a98d2cdfef213373a4
% Q3 figure source: a3970ef9eb51b87bd1c564a98d2cdfef213373a4

\section{模型检验}\label{sec:validation}

为检验结果可靠性，本文从独立输出验证、关键目标认证及候选池参数敏感性三个方面进行复核。三问正式方案均由独立程序重新读取交付CSV并重算约束与指标，避免仅依赖求解器内部状态。

\subsection{模型有效性与结果验证}\label{subsec:validation-results}

\subsubsection{独立验证与认证边界}

独立验证覆盖人员完整性、路线闭环、动态容量、燃油与加油等公共约束；问题三进一步检查人员时间窗、具体飞机兼容、任务冲突及30 min周转。三问正式结果及认证状态汇总于表~\ref{tab:validation-summary}。

\begin{table}[htbp]
\centering
\caption{三问正式结果验证与认证状态}
\label{tab:validation-summary}
\small
\begin{tabular}{c c c c}
\toprule
对象 & 服务结果 & 独立验证 & 认证说明\\
\midrule
问题一 & 1600人全部完成 & 通过 & 当前最好已验证可行解\\
问题二 & 4000人全部完成 & 通过 & 当前最好已验证可行解\\
问题三 & 3840+160人全部完成 & 通过 & 完整问题三未全局认证\\
问题三阶段B & 临时160/160 & 通过 & 受限模型内最优，相对间隙为0\\
\bottomrule
\end{tabular}
\end{table}

其中问题三阶段B在冻结的$B_0=35114$ min和当前有限时间可行架次列池下，临时人员目标的当前最好可行值与求解器最优界均为160，因此

\[
N_{\mathrm{temp},\Omega}^{*}=160.
\tag{9-1}\label{eq:q3-stageb-certification}
\]

160同时等于临时人员总数，故该目标达到人数意义上的理论绝对上限。该认证只针对阶段B的接纳目标；问题一、问题二以及问题三阶段A、阶段C仍按经验证可行结果解释，不扩大为完整路线空间的全局最优。

\subsection{稳定性与敏感性分析}\label{subsec:sensitivity}

\subsubsection{问题一候选加油参数敏感性}

该敏感性实验采用独立且固定的求解设置，仅用于比较不同加油候选保留数对结果的影响，并非对14759 min正式方案的重复求解。在其余设置保持一致的条件下，将加油候选保留数由4调整为8，结果见表~\ref{tab:q1-refuel-sensitivity}。

\begin{table}[htbp]
\centering
\caption{问题一候选加油参数敏感性}
\label{tab:q1-refuel-sensitivity}
\small
\setlength{\tabcolsep}{4pt}
\begin{tabular}{c c c c c c}
\toprule
加油候选保留数 & 路线池 & $T^{air}$/min & $T^{pax}$/人$\cdot$min & 燃油/kg & 架次数\\
\midrule
4 & 1603 & 14837 & 127784 & 119359.5 & 88\\
8 & 1622 & 14908 & 124166 & 119365.4 & 89\\
\bottomrule
\end{tabular}
\end{table}

由表~\ref{tab:q1-refuel-sensitivity}可见，参数由4增至8后，路线池仅扩大1.185\%，飞机时间变化0.479\%，燃油变化不足0.005\%。人员在途时间下降2.831\%，总体表明主要运行指标对该参数较稳定。不同指标的小幅反向变化来自有限时间词典序搜索，不能据此认为更大的候选池必然更差。

\subsubsection{问题三时间列密度敏感性}

问题三在基准采样方案“最早--中点--最晚”的基础上构造加密采样方案，并增加25\%和75\%分位点。图~\ref{fig:q3-time-density}显示，时间可行架次列与服务选项规模分别增加约85\%和94\%；与此同时，加密采样方案的阶段A及阶段C飞机使用时间均下降，燃油由301164.1 kg降至289768.1 kg，座位利用率由44.07\%升至45.52\%，说明更细时间嵌入能够提供更多有效排班组合。

\begin{figure}[htbp]
    \centering
    \includegraphics[width=0.96\textwidth]{figures/14时间列密度敏感性.pdf}
    \caption{问题三时间列密度敏感性分析}
    \label{fig:q3-time-density}
\end{figure}

但在相同限时求解条件下，阶段B当前可行接纳数由143降至112。由于加密采样候选集是在基准方案上扩展得到，可行域并未缩小，因此该差异应归因于模型规模增大后的组合搜索困难，而不能解释为“时间列越密，可服务人数越少”。这一判断还得到正式求解的验证：在更充分搜索后，阶段B最终达到160/160，并获得相对间隙为0的严格认证。因此143和112仅反映有限时间内当前最好可行解的质量，而不是临时人员服务能力上限。

综上，三问正式交付文件均通过独立验证；问题三阶段B的临时人员接纳目标得到严格认证。现有敏感性实验表明，问题一对加油候选保留数总体稳定，而问题三时间列加密能够提升排班灵活性，但会显著增加有限时间搜索负担。独立可行性复核、受限目标认证与参数扰动实验共同构成了本文的模型检验链。

% mcm-paper source binding (non-rendering metadata)
% Repository HEAD: 66f49838fe3060c9c0286cb5aa87613ac20b9fc3
% Section role: model evaluation and future improvements; no new numerical claims introduced.

\section{模型评价与改进}\label{sec:evaluation}

\subsection{模型优点}\label{subsec:model-strengths}

本文以统一架次可行性核为基础，逐步将单向出海、联合取送和有限机队多日排班纳入同一建模框架。问题一通过物理路线执行次数与人员流变量分离，避免将人员配置固化在候选列中；问题二进一步利用半开区间动态容量刻画同站座位即时复用；问题三则在此基础上引入时间可行架次列和具体飞机指派，使模型能够自然承接时间窗、周转和单机冲突约束。三问之间具有清晰的递进关系，公共规则可以复用，新增约束也具有明确物理含义。

在求解层面，候选路线生成与整数主问题相分离，使大规模组合空间能够通过“启发式构池+精确整数优化”的方式处理。路线的容量、燃油、加油和机场闭环等条件在进入主问题前即完成可行性筛选，从而减少无效变量；求解后再将聚合人员流回填为具体人员和实际架次，提高了模型结果向可执行运输方案转化的能力。

多目标处理采用词典序优化，而非人为设定加权系数，使飞机使用时间、人员在途时间、燃油、架次数和座位利用率之间的优先级保持清晰。特别地，问题三将非临时任务基准、临时任务接纳和次级指标细化分成A--C三个阶段，使“完成刚性任务”和“利用剩余能力提高临时任务满足率”在逻辑上相互分离。其中阶段B在固定时间预算和当前有限时间可行架次列池下实现160/160全部接纳，并获得相对间隙为0的严格认证，也增强了阶段性结论的可信度。

此外，三问最终交付方案均通过独立校验程序复核。验证过程脱离求解器内部对象，重新检查人员完整性、容量、燃油及路线等硬约束，问题三还进一步检查时间窗、具体飞机冲突和30 min周转；同时配合候选池与时间列密度敏感性实验，使结果的可行性、稳定性和认证边界均有对应证据支撑。

\subsection{模型不足与改进}\label{subsec:model-limitations}

首先，问题一、问题二以及问题三的阶段A和阶段C均建立在有限候选路线或时间可行架次列池上，候选集由启发式规则生成，因此当前结果不能等同于完整可行空间上的全局最优。后续可引入列生成或定价思想，根据主问题对偶信息动态产生更有价值的路线，并进一步结合Branch-and-Price或自适应扩池策略，减少人工限定候选池所带来的最优性缺口。

其次，问题三采用临界起飞时刻离散化以控制模型规模。时间列密度实验表明，更密的时间采样能够提供更丰富的排班组合，但也显著增加整数搜索负担。后续可采用自适应时间离散方法，仅在紧张时间窗或高冲突区域增密候选时刻；也可将7日排班分解为日级或滚动时域子问题，再通过跨日资源协调逐步修正，以提高大规模时空调度的求解效率。

再次，当前词典序优化的后级目标依赖上一级获得的当前解，其中问题三阶段C尚未完成全部五级最优性认证，座位利用率阶段也未形成新的正式改进结果。后续可增加后级目标求解预算，并结合$\varepsilon$-约束法或Pareto前沿分析，在保持主目标允许微小退让的条件下考察人员时间、燃油和利用率之间的替代关系，从而获得更丰富的决策信息。

最后，现有敏感性分析主要覆盖问题一的加油候选保留数和问题三的时间列密度，问题二尚未针对路线池构造参数开展独立扰动实验。后续可系统改变路线池规模、成对插入候选范围和求解时间预算，并采用多随机种子重复求解，以进一步检验联合取送结果的稳定性。若进一步考虑实际运行中的天气、需求变化、飞行时间波动或临时延误，可在现有确定性模型上增加时间缓冲、情景约束或随机/鲁棒优化机制。已有研究分别采用情景化样本平均近似处理带时间窗取送问题中的随机需求\cite{ghilas2016stochastic}，以及采用两阶段随机规划刻画异构飞行器随机速度与时间窗约束\cite{jia2018}，可为后续不确定性扩展提供参考。

\section*{AI 工具使用声明}
本参赛队在竞赛过程中使用了 AI 工具，主要用于题目分析、代码开发与调试、结果审查、模型与算法文档整理、图表与排版辅助、参考文献整理及论文表述，详细使用情况见支撑材料中的《AI 工具使用详情.pdf》。

% References are grounded in the user-provided literature package.
% Numbering follows the order of \bibitem entries and is cited in the main text via \cite{...}.

\renewcommand{\refname}{参考文献}
\begin{thebibliography}{99}
\setlength{\itemsep}{0.35em}
\setlength{\parsep}{0pt}

\bibitem{liang2024}
梁永宏, 葛显龙, 王博, 等. 异构车队多中心交付的低碳车辆路径问题与算法[J]. 重庆交通大学学报(自然科学版), 2024, 43(8): 60--68. DOI: 10.3969/j.issn.1674-0696.2024.08.08.

\bibitem{sethanan2020}
SETHANAN K, JAMRUS T. Hybrid differential evolution algorithm and genetic operator for multi-trip vehicle routing problem with backhauls and heterogeneous fleet in the beverage logistics industry[J]. Computers \& Industrial Engineering, 2020, 146: 106571. DOI: 10.1016/j.cie.2020.106571.

\bibitem{nafstad2021}
NAFSTAD G M, HAUGSETH A, HØYLAND V, STÅLHANE M. An exact solution method for a rich helicopter flight scheduling problem arising in offshore oil and gas logistics[J]. Computers \& Operations Research, 2021, 128: 105158. DOI: 10.1016/j.cor.2020.105158.

\bibitem{sundar2017}
SUNDAR K, VENKATACHALAM S, RATHINAM S. Analysis of mixed-integer linear programming formulations for a fuel-constrained multiple vehicle routing problem[J]. Unmanned Systems, 2017, 5(4): 197--207. DOI: 10.1142/S2301385017500091.

\bibitem{ghilas2016alns}
GHILAS V, DEMIR E, VAN WOENSEL T. An adaptive large neighborhood search heuristic for the pickup and delivery problem with time windows and scheduled lines[J]. Computers \& Operations Research, 2016, 72: 12--30. DOI: 10.1016/j.cor.2016.01.018.

\bibitem{ghilas2016stochastic}
GHILAS V, DEMIR E, VAN WOENSEL T. A scenario-based planning for the pickup and delivery problem with time windows, scheduled lines and stochastic demands[J]. Transportation Research Part B: Methodological, 2016, 91: 34--51. DOI: 10.1016/j.trb.2016.04.015.

\bibitem{jia2018}
JIA Z, YU J, AI X, et al. Cooperative multiple task assignment problem with stochastic velocities and time windows for heterogeneous unmanned aerial vehicles using a genetic algorithm[J]. Aerospace Science and Technology, 2018, 76: 112--125. DOI: 10.1016/j.ast.2018.01.025.

\end{thebibliography}

\clearpage
\section*{附录 A\quad 支撑材料与源程序说明}
\label{sec:appendix-support}

为便于复核与复现，本文将与三问建模、求解、模型检验和论文图表直接相关的程序、结果文件、自主查阅资料及 AI 使用说明随支撑材料一并提交。支撑材料与论文中的模型、数值和图表采用同一正式结果口径；赛题提供的原始数据不重复作为自主资料提交。

\subsection*{A.1\quad 支撑材料文件列表}

最终支撑材料压缩包按 \texttt{Support\_Materials/} 目录组织，实际文件如下。该清单与提交候选 ZIP 逐项对应。

\begin{longtable}{>{\raggedleft\arraybackslash}p{0.06\linewidth}>{\raggedright\arraybackslash}p{0.86\linewidth}}
\toprule
序号 & 相对路径 \\
\midrule
\endfirsthead
\toprule
序号 & 相对路径 \\
\midrule
\endhead
1 & \nolinkurl{Support_Materials/01_Programming/FILE_LIST.txt} \\
2 & \nolinkurl{Support_Materials/01_Programming/README.md} \\
3 & \nolinkurl{Support_Materials/01_Programming/manifest_sha256.txt} \\
4 & \nolinkurl{Support_Materials/01_Programming/plotting/q1/fig04_q1_history.m} \\
5 & \nolinkurl{Support_Materials/01_Programming/plotting/q1/fig05_q1_stages.m} \\
6 & \nolinkurl{Support_Materials/01_Programming/plotting/q1/fig06_q1_typical_flight.m} \\
7 & \nolinkurl{Support_Materials/01_Programming/plotting/q2/fig07_q2_joint_vs_baseline.m} \\
8 & \nolinkurl{Support_Materials/01_Programming/plotting/q2/fig08_q2_tpax_breakdown.m} \\
9 & \nolinkurl{Support_Materials/01_Programming/plotting/q2/fig09_q2_solver_stages.m} \\
10 & \nolinkurl{Support_Materials/01_Programming/plotting/q2/fig10_q2_dynamic_capacity.m} \\
11 & \nolinkurl{Support_Materials/01_Programming/plotting/q3/fig11_q3_pipeline.m} \\
12 & \nolinkurl{Support_Materials/01_Programming/plotting/q3/fig12_q3_fleet_gantt.m} \\
13 & \nolinkurl{Support_Materials/01_Programming/plotting/q3/fig13_q3_temporary_bottleneck.m} \\
14 & \nolinkurl{Support_Materials/01_Programming/plotting/q3/fig14_q3_time_density_sensitivity.m} \\
15 & \nolinkurl{Support_Materials/01_Programming/result_csv/q1-assignments.csv} \\
16 & \nolinkurl{Support_Materials/01_Programming/result_csv/q1-routes.csv} \\
17 & \nolinkurl{Support_Materials/01_Programming/result_csv/q2-assignments.csv} \\
18 & \nolinkurl{Support_Materials/01_Programming/result_csv/q2-routes.csv} \\
19 & \nolinkurl{Support_Materials/01_Programming/result_csv/q3-assignments.csv} \\
20 & \nolinkurl{Support_Materials/01_Programming/result_csv/q3-routes.csv} \\
21 & \nolinkurl{Support_Materials/01_Programming/source_code/common/__init__.py} \\
22 & \nolinkurl{Support_Materials/01_Programming/source_code/common/config.py} \\
23 & \nolinkurl{Support_Materials/01_Programming/source_code/common/data_loader.py} \\
24 & \nolinkurl{Support_Materials/01_Programming/source_code/common/feasibility.py} \\
25 & \nolinkurl{Support_Materials/01_Programming/source_code/common/grouping.py} \\
26 & \nolinkurl{Support_Materials/01_Programming/source_code/common/models.py} \\
27 & \nolinkurl{Support_Materials/01_Programming/source_code/common/paths.py} \\
28 & \nolinkurl{Support_Materials/01_Programming/source_code/data/official_data_placement.txt} \\
29 & \nolinkurl{Support_Materials/01_Programming/source_code/pyproject.toml} \\
30 & \nolinkurl{Support_Materials/01_Programming/source_code/q1/__init__.py} \\
31 & \nolinkurl{Support_Materials/01_Programming/source_code/q1/__main__.py} \\
32 & \nolinkurl{Support_Materials/01_Programming/source_code/q1/algorithms/__init__.py} \\
33 & \nolinkurl{Support_Materials/01_Programming/source_code/q1/algorithms/clusters.py} \\
34 & \nolinkurl{Support_Materials/01_Programming/source_code/q1/algorithms/flight_builder.py} \\
35 & \nolinkurl{Support_Materials/01_Programming/source_code/q1/algorithms/flow_master.py} \\
36 & \nolinkurl{Support_Materials/01_Programming/source_code/q1/algorithms/greedy.py} \\
37 & \nolinkurl{Support_Materials/01_Programming/source_code/q1/algorithms/route_dp.py} \\
38 & \nolinkurl{Support_Materials/01_Programming/source_code/q1/config.py} \\
39 & \nolinkurl{Support_Materials/01_Programming/source_code/q1/io/__init__.py} \\
40 & \nolinkurl{Support_Materials/01_Programming/source_code/q1/io/csv_validator.py} \\
41 & \nolinkurl{Support_Materials/01_Programming/source_code/q1/io/csv_writer.py} \\
42 & \nolinkurl{Support_Materials/01_Programming/source_code/q1/outputs/q1-assignments.csv} \\
43 & \nolinkurl{Support_Materials/01_Programming/source_code/q1/outputs/q1-history.csv} \\
44 & \nolinkurl{Support_Materials/01_Programming/source_code/q1/outputs/q1-routes.csv} \\
45 & \nolinkurl{Support_Materials/01_Programming/source_code/q1/outputs/q1-summary.json} \\
46 & \nolinkurl{Support_Materials/01_Programming/source_code/q1/outputs/sensitivity/sensitivity_q1.csv} \\
47 & \nolinkurl{Support_Materials/01_Programming/source_code/q1/run_q1.py} \\
48 & \nolinkurl{Support_Materials/01_Programming/source_code/q1/run_sensitivity.py} \\
49 & \nolinkurl{Support_Materials/01_Programming/source_code/q2/__init__.py} \\
50 & \nolinkurl{Support_Materials/01_Programming/source_code/q2/__main__.py} \\
51 & \nolinkurl{Support_Materials/01_Programming/source_code/q2/algorithms/__init__.py} \\
52 & \nolinkurl{Support_Materials/01_Programming/source_code/q2/algorithms/master.py} \\
53 & \nolinkurl{Support_Materials/01_Programming/source_code/q2/algorithms/realize.py} \\
54 & \nolinkurl{Support_Materials/01_Programming/source_code/q2/algorithms/route_pool.py} \\
55 & \nolinkurl{Support_Materials/01_Programming/source_code/q2/config.py} \\
56 & \nolinkurl{Support_Materials/01_Programming/source_code/q2/io/__init__.py} \\
57 & \nolinkurl{Support_Materials/01_Programming/source_code/q2/io/csv_validator.py} \\
58 & \nolinkurl{Support_Materials/01_Programming/source_code/q2/io/csv_writer.py} \\
59 & \nolinkurl{Support_Materials/01_Programming/source_code/q2/outputs/q2-assignments.csv} \\
60 & \nolinkurl{Support_Materials/01_Programming/source_code/q2/outputs/q2-detour-summary.csv} \\
61 & \nolinkurl{Support_Materials/01_Programming/source_code/q2/outputs/q2-history.csv} \\
62 & \nolinkurl{Support_Materials/01_Programming/source_code/q2/outputs/q2-routes.csv} \\
63 & \nolinkurl{Support_Materials/01_Programming/source_code/q2/outputs/q2-summary.json} \\
64 & \nolinkurl{Support_Materials/01_Programming/source_code/q2/outputs/q2-tpax-aggregate.csv} \\
65 & \nolinkurl{Support_Materials/01_Programming/source_code/q2/run_q2.py} \\
66 & \nolinkurl{Support_Materials/01_Programming/source_code/q3/__init__.py} \\
67 & \nolinkurl{Support_Materials/01_Programming/source_code/q3/__main__.py} \\
68 & \nolinkurl{Support_Materials/01_Programming/source_code/q3/algorithms/__init__.py} \\
69 & \nolinkurl{Support_Materials/01_Programming/source_code/q3/algorithms/lex.py} \\
70 & \nolinkurl{Support_Materials/01_Programming/source_code/q3/algorithms/master.py} \\
71 & \nolinkurl{Support_Materials/01_Programming/source_code/q3/algorithms/realize.py} \\
72 & \nolinkurl{Support_Materials/01_Programming/source_code/q3/algorithms/stage_a_pool.py} \\
73 & \nolinkurl{Support_Materials/01_Programming/source_code/q3/algorithms/timed_pool.py} \\
74 & \nolinkurl{Support_Materials/01_Programming/source_code/q3/analyze_temporary_bottleneck.py} \\
75 & \nolinkurl{Support_Materials/01_Programming/source_code/q3/config.py} \\
76 & \nolinkurl{Support_Materials/01_Programming/source_code/q3/fleet.py} \\
77 & \nolinkurl{Support_Materials/01_Programming/source_code/q3/io/__init__.py} \\
78 & \nolinkurl{Support_Materials/01_Programming/source_code/q3/io/csv_validator.py} \\
79 & \nolinkurl{Support_Materials/01_Programming/source_code/q3/io/csv_writer.py} \\
80 & \nolinkurl{Support_Materials/01_Programming/source_code/q3/outputs/formal_stage_a_summary.json} \\
81 & \nolinkurl{Support_Materials/01_Programming/source_code/q3/outputs/formal_stage_b_summary.json} \\
82 & \nolinkurl{Support_Materials/01_Programming/source_code/q3/outputs/q3-assignments.csv} \\
83 & \nolinkurl{Support_Materials/01_Programming/source_code/q3/outputs/q3-routes.csv} \\
84 & \nolinkurl{Support_Materials/01_Programming/source_code/q3/outputs/q3-summary.json} \\
85 & \nolinkurl{Support_Materials/01_Programming/source_code/q3/outputs/sensitivity/time_density_sensitivity.csv} \\
86 & \nolinkurl{Support_Materials/01_Programming/source_code/q3/outputs/temporary_bottleneck_summary.json} \\
87 & \nolinkurl{Support_Materials/01_Programming/source_code/q3/run_q3.py} \\
88 & \nolinkurl{Support_Materials/01_Programming/source_code/q3/run_time_density_sensitivity.py} \\
89 & \nolinkurl{Support_Materials/01_Programming/source_code/requirements.txt} \\
90 & \nolinkurl{Support_Materials/01_Programming/source_code/scripts/audit_tpax_breakdown.py} \\
91 & \nolinkurl{Support_Materials/01_Programming/source_code/scripts/verify_frozen_results.py} \\
92 & \nolinkurl{Support_Materials/02_Literature/REFERENCE_INDEX.md} \\
93 & \nolinkurl{Support_Materials/02_Literature/Ref01_Liang_2024_Heterogeneous_Multicenter_LowCarbon_VRP.pdf} \\
94 & \nolinkurl{Support_Materials/02_Literature/Ref02_Sethanan_Jamrus_2020_MultiTrip_Heterogeneous_Fleet.pdf} \\
95 & \nolinkurl{Support_Materials/02_Literature/Ref03_Nafstad_2021_Offshore_Helicopter_Scheduling.pdf} \\
96 & \nolinkurl{Support_Materials/02_Literature/Ref04_Sundar_2017_FuelConstrained_MultiVehicle_Routing.pdf} \\
97 & \nolinkurl{Support_Materials/02_Literature/Ref05_Ghilas_2016_ALNS_Pickup_Delivery_Scheduled_Lines.pdf} \\
98 & \nolinkurl{Support_Materials/02_Literature/Ref06_Ghilas_2016_Stochastic_Pickup_Delivery_Scheduled_Lines.pdf} \\
99 & \nolinkurl{Support_Materials/02_Literature/Ref07_Jia_2018_Stochastic_UAV_Task_Assignment.pdf} \\
100 & \nolinkurl{Support_Materials/03_Figures/FIGURE_INDEX.md} \\
101 & \nolinkurl{Support_Materials/03_Figures/Fig01_Unified_Framework.png} \\
102 & \nolinkurl{Support_Materials/03_Figures/Fig02_Network_Data_Overview.png} \\
103 & \nolinkurl{Support_Materials/03_Figures/Fig04_Q1_Solve_History.pdf} \\
104 & \nolinkurl{Support_Materials/03_Figures/Fig05_Q1_Lex_Status.pdf} \\
105 & \nolinkurl{Support_Materials/03_Figures/Fig06_Q1_Load_Utilization.pdf} \\
106 & \nolinkurl{Support_Materials/03_Figures/Fig07_Q2_Joint_vs_Baseline.pdf} \\
107 & \nolinkurl{Support_Materials/03_Figures/Fig08_Q2_Passenger_Time_Tradeoff.pdf} \\
108 & \nolinkurl{Support_Materials/03_Figures/Fig09_Q2_Lex_Status.pdf} \\
109 & \nolinkurl{Support_Materials/03_Figures/Fig10_Q2_Dynamic_Capacity.pdf} \\
110 & \nolinkurl{Support_Materials/03_Figures/Fig11_Q3_Stages.pdf} \\
111 & \nolinkurl{Support_Materials/03_Figures/Fig12_Q3_Fleet_Gantt.pdf} \\
112 & \nolinkurl{Support_Materials/03_Figures/Fig13_Q3_Temporary_Compatibility.pdf} \\
113 & \nolinkurl{Support_Materials/03_Figures/Fig14_Q3_Time_Density_Sensitivity.pdf} \\
114 & \texttt{Support\_Materials/}AI 工具使用详情.pdf \\
115 & \nolinkurl{Support_Materials/FILE_LIST.txt} \\
116 & \nolinkurl{Support_Materials/README.md} \\
117 & \nolinkurl{Support_Materials/manifest_sha256.txt} \\
\bottomrule
\end{longtable}

\subsection*{A.2\quad 源程序与复现说明}

三问的正式结果由对应求解程序在赛题数据上运行得到。问题一和问题二采用有限物理路线池与人员流整数规划，问题三进一步引入时间可行架次列、具体飞机指派与多日机队排班约束。支撑材料同时保存路线与人员分配结果、求解汇总、敏感性实验以及独立校验程序。

独立校验程序不依赖求解器内部变量，而是重新读取赛题数据与最终路线、人员分配等交付文件，对人员完整性、路线闭环、逐航段容量、燃油与加油规则进行核验；问题三额外核验人员时间窗、具体飞机兼容性、单机任务冲突和 30 min 周转约束。论文中的主要数值以经上述独立复核的正式交付方案为准。

程序运行环境、官方数据放置方式以及 Q1/Q2/Q3 复现命令见支撑材料编程目录中的 README.md。整数规划中的未认证阶段采用限时搜索，不同求解器版本、硬件和线程调度可能使重新运行获得不同的当前最好可行解，因此论文严格区分当前最好可行值、求解器最优界以及受限候选池内获得认证的最优值。

\subsection*{A.3\quad 完整源程序}

按照竞赛格式规范，下面给出建模、求解、独立校验、敏感性实验、结果审计以及论文图表生成所使用的完整源程序。为避免论文附录与支撑材料版本不一致，本附录从与最终提交候选包同源的 \texttt{appendix\_code/} 文件直接载入。

\subsubsection*{A.3.1\quad 公共模块与运行环境}

\paragraph{\nolinkurl{appendix_code/source_code/common/__init__.py}}

\VerbatimInput[fontsize=\tiny,breaklines=true,breakanywhere=true,numbers=left,numbersep=3pt,tabsize=4,frame=single]{appendix_code/source_code/common/__init__.py}

\paragraph{\nolinkurl{appendix_code/source_code/common/config.py}}

\VerbatimInput[fontsize=\tiny,breaklines=true,breakanywhere=true,numbers=left,numbersep=3pt,tabsize=4,frame=single]{appendix_code/source_code/common/config.py}

\paragraph{\nolinkurl{appendix_code/source_code/common/data_loader.py}}

\VerbatimInput[fontsize=\tiny,breaklines=true,breakanywhere=true,numbers=left,numbersep=3pt,tabsize=4,frame=single]{appendix_code/source_code/common/data_loader.py}

\paragraph{\nolinkurl{appendix_code/source_code/common/feasibility.py}}

\VerbatimInput[fontsize=\tiny,breaklines=true,breakanywhere=true,numbers=left,numbersep=3pt,tabsize=4,frame=single]{appendix_code/source_code/common/feasibility.py}

\paragraph{\nolinkurl{appendix_code/source_code/common/grouping.py}}

\VerbatimInput[fontsize=\tiny,breaklines=true,breakanywhere=true,numbers=left,numbersep=3pt,tabsize=4,frame=single]{appendix_code/source_code/common/grouping.py}

\paragraph{\nolinkurl{appendix_code/source_code/common/models.py}}

\VerbatimInput[fontsize=\tiny,breaklines=true,breakanywhere=true,numbers=left,numbersep=3pt,tabsize=4,frame=single]{appendix_code/source_code/common/models.py}

\paragraph{\nolinkurl{appendix_code/source_code/common/paths.py}}

\VerbatimInput[fontsize=\tiny,breaklines=true,breakanywhere=true,numbers=left,numbersep=3pt,tabsize=4,frame=single]{appendix_code/source_code/common/paths.py}

\paragraph{\nolinkurl{appendix_code/source_code/data/official_data_placement.txt}}

\VerbatimInput[fontsize=\tiny,breaklines=true,breakanywhere=true,numbers=left,numbersep=3pt,tabsize=4,frame=single]{appendix_code/source_code/data/official_data_placement.txt}

\paragraph{\nolinkurl{appendix_code/source_code/pyproject.toml}}

\VerbatimInput[fontsize=\tiny,breaklines=true,breakanywhere=true,numbers=left,numbersep=3pt,tabsize=4,frame=single]{appendix_code/source_code/pyproject.toml}

\paragraph{\nolinkurl{appendix_code/source_code/requirements.txt}}

\VerbatimInput[fontsize=\tiny,breaklines=true,breakanywhere=true,numbers=left,numbersep=3pt,tabsize=4,frame=single]{appendix_code/source_code/requirements.txt}

\subsubsection*{A.3.2\quad 问题一源程序}

\paragraph{\nolinkurl{appendix_code/source_code/q1/__init__.py}}

\VerbatimInput[fontsize=\tiny,breaklines=true,breakanywhere=true,numbers=left,numbersep=3pt,tabsize=4,frame=single]{appendix_code/source_code/q1/__init__.py}

\paragraph{\nolinkurl{appendix_code/source_code/q1/__main__.py}}

\VerbatimInput[fontsize=\tiny,breaklines=true,breakanywhere=true,numbers=left,numbersep=3pt,tabsize=4,frame=single]{appendix_code/source_code/q1/__main__.py}

\paragraph{\nolinkurl{appendix_code/source_code/q1/algorithms/__init__.py}}

\VerbatimInput[fontsize=\tiny,breaklines=true,breakanywhere=true,numbers=left,numbersep=3pt,tabsize=4,frame=single]{appendix_code/source_code/q1/algorithms/__init__.py}

\paragraph{\nolinkurl{appendix_code/source_code/q1/algorithms/clusters.py}}

\VerbatimInput[fontsize=\tiny,breaklines=true,breakanywhere=true,numbers=left,numbersep=3pt,tabsize=4,frame=single]{appendix_code/source_code/q1/algorithms/clusters.py}

\paragraph{\nolinkurl{appendix_code/source_code/q1/algorithms/flight_builder.py}}

\VerbatimInput[fontsize=\tiny,breaklines=true,breakanywhere=true,numbers=left,numbersep=3pt,tabsize=4,frame=single]{appendix_code/source_code/q1/algorithms/flight_builder.py}

\paragraph{\nolinkurl{appendix_code/source_code/q1/algorithms/flow_master.py}}

\VerbatimInput[fontsize=\tiny,breaklines=true,breakanywhere=true,numbers=left,numbersep=3pt,tabsize=4,frame=single]{appendix_code/source_code/q1/algorithms/flow_master.py}

\paragraph{\nolinkurl{appendix_code/source_code/q1/algorithms/greedy.py}}

\VerbatimInput[fontsize=\tiny,breaklines=true,breakanywhere=true,numbers=left,numbersep=3pt,tabsize=4,frame=single]{appendix_code/source_code/q1/algorithms/greedy.py}

\paragraph{\nolinkurl{appendix_code/source_code/q1/algorithms/route_dp.py}}

\VerbatimInput[fontsize=\tiny,breaklines=true,breakanywhere=true,numbers=left,numbersep=3pt,tabsize=4,frame=single]{appendix_code/source_code/q1/algorithms/route_dp.py}

\paragraph{\nolinkurl{appendix_code/source_code/q1/config.py}}

\VerbatimInput[fontsize=\tiny,breaklines=true,breakanywhere=true,numbers=left,numbersep=3pt,tabsize=4,frame=single]{appendix_code/source_code/q1/config.py}

\paragraph{\nolinkurl{appendix_code/source_code/q1/io/__init__.py}}

\VerbatimInput[fontsize=\tiny,breaklines=true,breakanywhere=true,numbers=left,numbersep=3pt,tabsize=4,frame=single]{appendix_code/source_code/q1/io/__init__.py}

\paragraph{\nolinkurl{appendix_code/source_code/q1/io/csv_validator.py}}

\VerbatimInput[fontsize=\tiny,breaklines=true,breakanywhere=true,numbers=left,numbersep=3pt,tabsize=4,frame=single]{appendix_code/source_code/q1/io/csv_validator.py}

\paragraph{\nolinkurl{appendix_code/source_code/q1/io/csv_writer.py}}

\VerbatimInput[fontsize=\tiny,breaklines=true,breakanywhere=true,numbers=left,numbersep=3pt,tabsize=4,frame=single]{appendix_code/source_code/q1/io/csv_writer.py}

\paragraph{\nolinkurl{appendix_code/source_code/q1/run_q1.py}}

\VerbatimInput[fontsize=\tiny,breaklines=true,breakanywhere=true,numbers=left,numbersep=3pt,tabsize=4,frame=single]{appendix_code/source_code/q1/run_q1.py}

\paragraph{\nolinkurl{appendix_code/source_code/q1/run_sensitivity.py}}

\VerbatimInput[fontsize=\tiny,breaklines=true,breakanywhere=true,numbers=left,numbersep=3pt,tabsize=4,frame=single]{appendix_code/source_code/q1/run_sensitivity.py}

\subsubsection*{A.3.3\quad 问题二源程序}

\paragraph{\nolinkurl{appendix_code/source_code/q2/__init__.py}}

\VerbatimInput[fontsize=\tiny,breaklines=true,breakanywhere=true,numbers=left,numbersep=3pt,tabsize=4,frame=single]{appendix_code/source_code/q2/__init__.py}

\paragraph{\nolinkurl{appendix_code/source_code/q2/__main__.py}}

\VerbatimInput[fontsize=\tiny,breaklines=true,breakanywhere=true,numbers=left,numbersep=3pt,tabsize=4,frame=single]{appendix_code/source_code/q2/__main__.py}

\paragraph{\nolinkurl{appendix_code/source_code/q2/algorithms/__init__.py}}

\VerbatimInput[fontsize=\tiny,breaklines=true,breakanywhere=true,numbers=left,numbersep=3pt,tabsize=4,frame=single]{appendix_code/source_code/q2/algorithms/__init__.py}

\paragraph{\nolinkurl{appendix_code/source_code/q2/algorithms/master.py}}

\VerbatimInput[fontsize=\tiny,breaklines=true,breakanywhere=true,numbers=left,numbersep=3pt,tabsize=4,frame=single]{appendix_code/source_code/q2/algorithms/master.py}

\paragraph{\nolinkurl{appendix_code/source_code/q2/algorithms/realize.py}}

\VerbatimInput[fontsize=\tiny,breaklines=true,breakanywhere=true,numbers=left,numbersep=3pt,tabsize=4,frame=single]{appendix_code/source_code/q2/algorithms/realize.py}

\paragraph{\nolinkurl{appendix_code/source_code/q2/algorithms/route_pool.py}}

\VerbatimInput[fontsize=\tiny,breaklines=true,breakanywhere=true,numbers=left,numbersep=3pt,tabsize=4,frame=single]{appendix_code/source_code/q2/algorithms/route_pool.py}

\paragraph{\nolinkurl{appendix_code/source_code/q2/config.py}}

\VerbatimInput[fontsize=\tiny,breaklines=true,breakanywhere=true,numbers=left,numbersep=3pt,tabsize=4,frame=single]{appendix_code/source_code/q2/config.py}

\paragraph{\nolinkurl{appendix_code/source_code/q2/io/__init__.py}}

\VerbatimInput[fontsize=\tiny,breaklines=true,breakanywhere=true,numbers=left,numbersep=3pt,tabsize=4,frame=single]{appendix_code/source_code/q2/io/__init__.py}

\paragraph{\nolinkurl{appendix_code/source_code/q2/io/csv_validator.py}}

\VerbatimInput[fontsize=\tiny,breaklines=true,breakanywhere=true,numbers=left,numbersep=3pt,tabsize=4,frame=single]{appendix_code/source_code/q2/io/csv_validator.py}

\paragraph{\nolinkurl{appendix_code/source_code/q2/io/csv_writer.py}}

\VerbatimInput[fontsize=\tiny,breaklines=true,breakanywhere=true,numbers=left,numbersep=3pt,tabsize=4,frame=single]{appendix_code/source_code/q2/io/csv_writer.py}

\paragraph{\nolinkurl{appendix_code/source_code/q2/run_q2.py}}

\VerbatimInput[fontsize=\tiny,breaklines=true,breakanywhere=true,numbers=left,numbersep=3pt,tabsize=4,frame=single]{appendix_code/source_code/q2/run_q2.py}

\subsubsection*{A.3.4\quad 问题三源程序}

\paragraph{\nolinkurl{appendix_code/source_code/q3/__init__.py}}

\VerbatimInput[fontsize=\tiny,breaklines=true,breakanywhere=true,numbers=left,numbersep=3pt,tabsize=4,frame=single]{appendix_code/source_code/q3/__init__.py}

\paragraph{\nolinkurl{appendix_code/source_code/q3/__main__.py}}

\VerbatimInput[fontsize=\tiny,breaklines=true,breakanywhere=true,numbers=left,numbersep=3pt,tabsize=4,frame=single]{appendix_code/source_code/q3/__main__.py}

\paragraph{\nolinkurl{appendix_code/source_code/q3/algorithms/__init__.py}}

\VerbatimInput[fontsize=\tiny,breaklines=true,breakanywhere=true,numbers=left,numbersep=3pt,tabsize=4,frame=single]{appendix_code/source_code/q3/algorithms/__init__.py}

\paragraph{\nolinkurl{appendix_code/source_code/q3/algorithms/lex.py}}

\VerbatimInput[fontsize=\tiny,breaklines=true,breakanywhere=true,numbers=left,numbersep=3pt,tabsize=4,frame=single]{appendix_code/source_code/q3/algorithms/lex.py}

\paragraph{\nolinkurl{appendix_code/source_code/q3/algorithms/master.py}}

\VerbatimInput[fontsize=\tiny,breaklines=true,breakanywhere=true,numbers=left,numbersep=3pt,tabsize=4,frame=single]{appendix_code/source_code/q3/algorithms/master.py}

\paragraph{\nolinkurl{appendix_code/source_code/q3/algorithms/realize.py}}

\VerbatimInput[fontsize=\tiny,breaklines=true,breakanywhere=true,numbers=left,numbersep=3pt,tabsize=4,frame=single]{appendix_code/source_code/q3/algorithms/realize.py}

\paragraph{\nolinkurl{appendix_code/source_code/q3/algorithms/stage_a_pool.py}}

\VerbatimInput[fontsize=\tiny,breaklines=true,breakanywhere=true,numbers=left,numbersep=3pt,tabsize=4,frame=single]{appendix_code/source_code/q3/algorithms/stage_a_pool.py}

\paragraph{\nolinkurl{appendix_code/source_code/q3/algorithms/timed_pool.py}}

\VerbatimInput[fontsize=\tiny,breaklines=true,breakanywhere=true,numbers=left,numbersep=3pt,tabsize=4,frame=single]{appendix_code/source_code/q3/algorithms/timed_pool.py}

\paragraph{\nolinkurl{appendix_code/source_code/q3/analyze_temporary_bottleneck.py}}

\VerbatimInput[fontsize=\tiny,breaklines=true,breakanywhere=true,numbers=left,numbersep=3pt,tabsize=4,frame=single]{appendix_code/source_code/q3/analyze_temporary_bottleneck.py}

\paragraph{\nolinkurl{appendix_code/source_code/q3/config.py}}

\VerbatimInput[fontsize=\tiny,breaklines=true,breakanywhere=true,numbers=left,numbersep=3pt,tabsize=4,frame=single]{appendix_code/source_code/q3/config.py}

\paragraph{\nolinkurl{appendix_code/source_code/q3/fleet.py}}

\VerbatimInput[fontsize=\tiny,breaklines=true,breakanywhere=true,numbers=left,numbersep=3pt,tabsize=4,frame=single]{appendix_code/source_code/q3/fleet.py}

\paragraph{\nolinkurl{appendix_code/source_code/q3/io/__init__.py}}

\VerbatimInput[fontsize=\tiny,breaklines=true,breakanywhere=true,numbers=left,numbersep=3pt,tabsize=4,frame=single]{appendix_code/source_code/q3/io/__init__.py}

\paragraph{\nolinkurl{appendix_code/source_code/q3/io/csv_validator.py}}

\VerbatimInput[fontsize=\tiny,breaklines=true,breakanywhere=true,numbers=left,numbersep=3pt,tabsize=4,frame=single]{appendix_code/source_code/q3/io/csv_validator.py}

\paragraph{\nolinkurl{appendix_code/source_code/q3/io/csv_writer.py}}

\VerbatimInput[fontsize=\tiny,breaklines=true,breakanywhere=true,numbers=left,numbersep=3pt,tabsize=4,frame=single]{appendix_code/source_code/q3/io/csv_writer.py}

\paragraph{\nolinkurl{appendix_code/source_code/q3/run_q3.py}}

\VerbatimInput[fontsize=\tiny,breaklines=true,breakanywhere=true,numbers=left,numbersep=3pt,tabsize=4,frame=single]{appendix_code/source_code/q3/run_q3.py}

\paragraph{\nolinkurl{appendix_code/source_code/q3/run_time_density_sensitivity.py}}

\VerbatimInput[fontsize=\tiny,breaklines=true,breakanywhere=true,numbers=left,numbersep=3pt,tabsize=4,frame=single]{appendix_code/source_code/q3/run_time_density_sensitivity.py}

\subsubsection*{A.3.5\quad 结果审计与冻结结果核验程序}

\paragraph{\nolinkurl{appendix_code/source_code/scripts/audit_tpax_breakdown.py}}

\VerbatimInput[fontsize=\tiny,breaklines=true,breakanywhere=true,numbers=left,numbersep=3pt,tabsize=4,frame=single]{appendix_code/source_code/scripts/audit_tpax_breakdown.py}

\paragraph{\nolinkurl{appendix_code/source_code/scripts/verify_frozen_results.py}}

\VerbatimInput[fontsize=\tiny,breaklines=true,breakanywhere=true,numbers=left,numbersep=3pt,tabsize=4,frame=single]{appendix_code/source_code/scripts/verify_frozen_results.py}

\subsubsection*{A.3.6\quad 论文图表 MATLAB 程序}

\paragraph{\nolinkurl{appendix_code/plotting/q1/fig04_q1_history.m}}

\VerbatimInput[fontsize=\tiny,breaklines=true,breakanywhere=true,numbers=left,numbersep=3pt,tabsize=4,frame=single]{appendix_code/plotting/q1/fig04_q1_history.m}

\paragraph{\nolinkurl{appendix_code/plotting/q1/fig05_q1_stages.m}}

\VerbatimInput[fontsize=\tiny,breaklines=true,breakanywhere=true,numbers=left,numbersep=3pt,tabsize=4,frame=single]{appendix_code/plotting/q1/fig05_q1_stages.m}

\paragraph{\nolinkurl{appendix_code/plotting/q1/fig06_q1_typical_flight.m}}

\VerbatimInput[fontsize=\tiny,breaklines=true,breakanywhere=true,numbers=left,numbersep=3pt,tabsize=4,frame=single]{appendix_code/plotting/q1/fig06_q1_typical_flight.m}

\paragraph{\nolinkurl{appendix_code/plotting/q2/fig07_q2_joint_vs_baseline.m}}

\VerbatimInput[fontsize=\tiny,breaklines=true,breakanywhere=true,numbers=left,numbersep=3pt,tabsize=4,frame=single]{appendix_code/plotting/q2/fig07_q2_joint_vs_baseline.m}

\paragraph{\nolinkurl{appendix_code/plotting/q2/fig08_q2_tpax_breakdown.m}}

\VerbatimInput[fontsize=\tiny,breaklines=true,breakanywhere=true,numbers=left,numbersep=3pt,tabsize=4,frame=single]{appendix_code/plotting/q2/fig08_q2_tpax_breakdown.m}

\paragraph{\nolinkurl{appendix_code/plotting/q2/fig09_q2_solver_stages.m}}

\VerbatimInput[fontsize=\tiny,breaklines=true,breakanywhere=true,numbers=left,numbersep=3pt,tabsize=4,frame=single]{appendix_code/plotting/q2/fig09_q2_solver_stages.m}

\paragraph{\nolinkurl{appendix_code/plotting/q2/fig10_q2_dynamic_capacity.m}}

\VerbatimInput[fontsize=\tiny,breaklines=true,breakanywhere=true,numbers=left,numbersep=3pt,tabsize=4,frame=single]{appendix_code/plotting/q2/fig10_q2_dynamic_capacity.m}

\paragraph{\nolinkurl{appendix_code/plotting/q3/fig11_q3_pipeline.m}}

\VerbatimInput[fontsize=\tiny,breaklines=true,breakanywhere=true,numbers=left,numbersep=3pt,tabsize=4,frame=single]{appendix_code/plotting/q3/fig11_q3_pipeline.m}

\paragraph{\nolinkurl{appendix_code/plotting/q3/fig12_q3_fleet_gantt.m}}

\VerbatimInput[fontsize=\tiny,breaklines=true,breakanywhere=true,numbers=left,numbersep=3pt,tabsize=4,frame=single]{appendix_code/plotting/q3/fig12_q3_fleet_gantt.m}

\paragraph{\nolinkurl{appendix_code/plotting/q3/fig13_q3_temporary_bottleneck.m}}

\VerbatimInput[fontsize=\tiny,breaklines=true,breakanywhere=true,numbers=left,numbersep=3pt,tabsize=4,frame=single]{appendix_code/plotting/q3/fig13_q3_temporary_bottleneck.m}

\paragraph{\nolinkurl{appendix_code/plotting/q3/fig14_q3_time_density_sensitivity.m}}

\VerbatimInput[fontsize=\tiny,breaklines=true,breakanywhere=true,numbers=left,numbersep=3pt,tabsize=4,frame=single]{appendix_code/plotting/q3/fig14_q3_time_density_sensitivity.m}

\subsection*{A.4\quad AI 工具使用声明}

本参赛队在竞赛过程中使用了 AI 工具，主要用于题目分析、代码开发与调试、结果审查、模型与算法文档整理、图表与排版辅助、参考文献整理及论文表述，详细使用情况见支撑材料中的《AI 工具使用详情.pdf》。所使用的主要工具包括 DeepSeek（deepseekv4flash）、Claude Code（Anthropic，竞赛期间客户端当前最新模型）、ChatGPT（OpenAI，GPT-5.6 Sol）和 Codex（OpenAI，GPT-5.6 Sol）。所有进入论文的模型、程序修改、数值、图表与文字均经过队员人工审查；关键数值以正式程序输出和独立校验结果为依据。


\end{document}
